\documentclass[11pt]{article}
% ===========================================================================
%  Shared preamble for all MPM2D worksheets — input right after \documentclass.
%  (Files beginning with "_" are skipped by mpm2d-worksheet-publish.mjs.)
% ===========================================================================
\usepackage[letterpaper,top=2.2cm,bottom=1.9cm,left=1.6cm,right=1.6cm,headheight=28pt,footskip=24pt]{geometry}
\usepackage{amsmath,amssymb}
\usepackage[T1]{fontenc}
\usepackage[utf8]{inputenc}
\usepackage{lmodern}
\usepackage{xcolor}
\usepackage[most]{tcolorbox}
\usepackage{enumitem}
\usepackage{multicol}
\usepackage{fancyhdr}
\usepackage{tikz}
\usetikzlibrary{calc}
\usepackage{pgfplots}
\pgfplotsset{compat=1.18}

\definecolor{exblue}{HTML}{4A90E2}
\definecolor{exbg}{HTML}{E6F3FF}
\definecolor{qorange}{HTML}{E69138}
\definecolor{qbg}{HTML}{FFF7CC}
\definecolor{keygray}{HTML}{475569}
\definecolor{keybg}{HTML}{F1F5F9}
\definecolor{iagreen}{HTML}{14653B}
\definecolor{titleblue}{HTML}{1D4ED8}
% Rotating palette so each worked example gets its own colour.
\definecolor{ex0}{HTML}{2563A0}\definecolor{ex1}{HTML}{3B7D3B}\definecolor{ex2}{HTML}{8A5A00}
\definecolor{ex3}{HTML}{A3327A}\definecolor{ex4}{HTML}{6D28A3}\definecolor{ex5}{HTML}{B08900}
\definecolor{ex6}{HTML}{0E7490}\definecolor{ex7}{HTML}{9A3412}\definecolor{ex8}{HTML}{15803D}

\pagestyle{fancy}
\fancyhf{}
\fancyhead[L]{\footnotesize NAME: \underline{\hspace{3.4cm}}}
\fancyhead[C]{\textbf{Integration Academy}\\[-2pt]{\scriptsize Math Simplified \textperiodcentered{} Success Amplified}}
\fancyhead[R]{\footnotesize MPM2D}
\fancyfoot[L]{\scriptsize dr.merie@IntegrationAcademy.ca}
\fancyfoot[C]{\scriptsize\thepage}
\fancyfoot[R]{\scriptsize IntegrationAcademy.ca}
\renewcommand{\headrulewidth}{1.2pt}
\renewcommand{\footrulewidth}{0.4pt}
\setlength{\parindent}{0pt}
\setlength{\parskip}{4pt}

% Examples are NOT breakable, so each example + its solution stays on one page.
\newtcolorbox{example}[2]{colframe=#1,colback=#1!7,coltitle=white,colbacktitle=#1,fonttitle=\bfseries,title={#2},arc=3pt,boxrule=1pt}
\newtcolorbox{qbox}[1]{colframe=qorange,colback=qbg,coltitle=white,colbacktitle=qorange,fonttitle=\bfseries,title={#1},arc=3pt,breakable,boxrule=1pt}
\newtcolorbox{keybox}{colframe=keygray,colback=keybg,coltitle=white,colbacktitle=keygray,fonttitle=\bfseries,title={$\checkmark$\ Answer Key},arc=3pt,breakable,boxrule=1pt}
\newtcolorbox{recapbox}{colframe=keygray,colback=keybg,coltitle=white,colbacktitle=keygray,fonttitle=\bfseries,title={Question Recap},arc=3pt,breakable,boxrule=1pt}
\newtcolorbox{ideabox}{colframe=titleblue,colback=titleblue!6,coltitle=white,colbacktitle=titleblue,fonttitle=\bfseries,title={The Big Ideas},arc=3pt,breakable,boxrule=1.2pt}
\newtcolorbox{learnbox}{colframe=iagreen,colback=iagreen!5,coltitle=white,colbacktitle=iagreen,fonttitle=\bfseries,title={Concept Essentials},arc=3pt,breakable,boxrule=1.2pt}
\newcommand{\lh}[1]{\par\smallskip{\bfseries\color{iagreen}#1}\par\nobreak\smallskip}

\newcommand{\soln}{\par\smallskip\textit{\textbf{Solution.}}\ }
\newcommand{\workspace}[1]{\par\vspace{3pt}\fcolorbox{black!30}{white}{\begin{minipage}[t][#1][t]{\dimexpr\linewidth-2\fboxsep\relax}\ \end{minipage}}\par}
\newcommand{\sech}[1]{\par\vspace{8pt}{\Large\bfseries\color{iagreen}#1}\par\nobreak\vspace{2pt}{\color{black!20}\hrule height1pt}\vspace{6pt}}

% A compact coordinate plot sized for an example box: \eplot{xmin}{xmax}{ymin}{ymax}{body}
\newcommand{\eplot}[5]{\begin{center}\begin{tikzpicture}
\begin{axis}[width=6.8cm,height=4.4cm,axis lines=middle,xlabel={\tiny$x$},ylabel={\tiny$y$},
grid=both,grid style={gray!16},xmin=#1,xmax=#2,ymin=#3,ymax=#4,
xtick distance=2,ytick distance=2,tick label style={font=\tiny},axis line style={-},samples=120]
#5
\end{axis}\end{tikzpicture}\end{center}}

% A sine-curve plot (degrees on x, 0..360): \splot{ymin}{ymax}{body}
\newcommand{\splot}[3]{\begin{center}\begin{tikzpicture}
\begin{axis}[width=8.6cm,height=4.6cm,axis lines=middle,xlabel={\tiny$x\,(^\circ)$},ylabel={\tiny$y$},
grid=both,grid style={gray!16},xmin=0,xmax=360,ymin=#1,ymax=#2,
xtick distance=90,ytick distance=1,tick label style={font=\tiny},axis line style={-},samples=180]
#3
\end{axis}\end{tikzpicture}\end{center}}

% A small coordinate plot: \plot{xmin}{xmax}{ymin}{ymax}{body}
\newcommand{\plot}[5]{\begin{center}\begin{tikzpicture}
\begin{axis}[width=8.4cm,height=6cm,axis lines=middle,xlabel={\scriptsize$x$},ylabel={\scriptsize$y$},
grid=both,grid style={gray!16},xmin=#1,xmax=#2,ymin=#3,ymax=#4,
xtick distance=2,ytick distance=2,tick label style={font=\tiny},axis line style={-}]
#5
\end{axis}\end{tikzpicture}\end{center}}

% Right triangle (angle theta at bottom-left, right angle at bottom-right):
%   \rtri{drawBase}{drawHeight}{oppLabel}{adjLabel}{hypLabel}
\newcommand{\rtri}[5]{\begin{center}\begin{tikzpicture}[scale=0.85]
\coordinate (A) at (0,0);\coordinate (B) at (#1,0);\coordinate (C) at (#1,#2);
\draw[thick,exblue] (A)--(B)--(C)--cycle;
\draw[black] ($(B)+(-0.32,0)$)--($(B)+(-0.32,0.32)$)--($(B)+(0,0.32)$);
\node[below] at ($(A)!0.5!(B)$){#4};
\node[right] at ($(B)!0.5!(C)$){#3};
\node[above left] at ($(A)!0.5!(C)$){#5};
\node at (0.7,0.22){$\theta$};
\end{tikzpicture}\end{center}}

% General (oblique) triangle with vertex labels and opposite-side labels:
%   \gtri{Alabel}{Blabel}{Clabel}{aSide}{bSide}{cSide}
\newcommand{\gtri}[6]{\begin{center}\begin{tikzpicture}[scale=0.85]
\coordinate (A) at (0,0);\coordinate (B) at (5,0);\coordinate (C) at (1.6,3);
\draw[thick,exblue] (A)--(B)--(C)--cycle;
\node[below left] at (A){#1};\node[below right] at (B){#2};\node[above] at (C){#3};
\node[below] at ($(A)!0.5!(B)$){#6};
\node[right] at ($(B)!0.5!(C)$){#4};
\node[left] at ($(A)!0.5!(C)$){#5};
\end{tikzpicture}\end{center}}

\fancyhead[R]{\footnotesize MCR3U \textperiodcentered{} 2.3}
\begin{document}

\begin{center}
{\LARGE\bfseries\color{titleblue}Worksheet 2.3 --- Simplifying Rational Expressions}\\[2pt]
{\small Grade 11 Functions, University (MCR3U) \textperiodcentered{} Unit 2: Equivalent Algebraic Expressions}
\end{center}

\begin{ideabox}
Treat algebraic fractions like number fractions: factor, cancel common factors, and state restrictions.
\begin{itemize}[leftmargin=*,itemsep=2pt]
  \item Factor numerator and denominator, then cancel common \textbf{factors} (never terms in a sum).
  \item State restrictions: any value making an original denominator $0$.
  \item Multiply: cancel across, then multiply. Divide: multiply by the reciprocal.
\end{itemize}
\end{ideabox}

\begin{learnbox}
\lh{Fractions with polynomials}A \textbf{rational expression} is a ratio of polynomials. Simplify by factoring top and bottom and cancelling common factors.
\lh{State restrictions}Any value making the original denominator $0$ is excluded — list these before cancelling, since cancelling can hide them.
\lh{Only factors cancel}You may cancel matching \emph{factors}, never individual terms: $\dfrac{x+2}{x}$ does not become $2$.
\end{learnbox}

\sech{Worked Examples}

\begin{example}{ex0}{Example 1 --- Cancel a monomial}
Simplify $\dfrac{3x^2}{6x}$.\soln Divide coefficients and subtract exponents: $\dfrac{3}{6}x^{2-1}=\dfrac{x}{2},\ x\ne0$.
\end{example}

\begin{example}{ex1}{Example 2 --- Difference of squares}
Simplify $\dfrac{x^2-9}{x-3}$.\soln Factor the top: $\dfrac{(x-3)(x+3)}{x-3}$. \\ Cancel $(x-3)$: $x+3,\ x\ne3$.
\end{example}

\begin{example}{ex2}{Example 3 --- Trinomial over binomial}
Simplify $\dfrac{x^2+5x+6}{x+2}$.\soln $\dfrac{(x+2)(x+3)}{x+2}=x+3,\ x\ne-2$.
\end{example}

\begin{example}{ex3}{Example 4 --- Trinomial over trinomial}
Simplify $\dfrac{x^2-x-6}{x^2-9}$.\soln Factor both: $\dfrac{(x-3)(x+2)}{(x-3)(x+3)}$. \\ Cancel $(x-3)$: $\dfrac{x+2}{x+3},\ x\ne\pm3$.
\end{example}

\begin{example}{ex4}{Example 5 --- Common factor}
Simplify $\dfrac{2x+4}{x^2+2x}$.\soln $\dfrac{2(x+2)}{x(x+2)}=\dfrac{2}{x},\ x\ne0,-2$.
\end{example}

\begin{example}{ex5}{Example 6 --- Multiply}
Simplify $\dfrac{x^2-4}{x}\cdot\dfrac{3}{x+2}$.\soln Factor and cancel across: $\dfrac{(x-2)(x+2)}{x}\cdot\dfrac{3}{x+2}=\dfrac{3(x-2)}{x},\ x\ne0,-2$.
\end{example}

\begin{example}{ex6}{Example 7 --- Multiply to 1}
Simplify $\dfrac{x+1}{x-2}\cdot\dfrac{x-2}{x+1}$.\soln Everything cancels: $1,\ x\ne2,-1$.
\end{example}

\begin{example}{ex7}{Example 8 --- Divide}
Simplify $\dfrac{x}{x+3}\div\dfrac{x^2}{x+3}$.\soln Multiply by the reciprocal: $\dfrac{x}{x+3}\cdot\dfrac{x+3}{x^2}=\dfrac{1}{x},\ x\ne-3,0$.
\end{example}

\begin{example}{ex8}{Example 9 --- Divide with factoring}
Simplify $\dfrac{x^2-1}{x}\div\dfrac{x+1}{x}$.\soln $\dfrac{x^2-1}{x}\cdot\dfrac{x}{x+1}=\dfrac{(x-1)(x+1)}{x+1}=x-1,\ x\ne0,-1$.
\end{example}

\clearpage
\sech{Your Turn --- Show Your Work}
For each question, show all steps clearly in the space provided.

\begin{qbox}{Question 1}
Simplify $\dfrac{4x^3}{8x}$ and state restrictions.\workspace{2.4cm}
\end{qbox}

\begin{qbox}{Question 2}
Simplify $\dfrac{x^2-25}{x+5}$.\workspace{2.4cm}
\end{qbox}

\begin{qbox}{Question 3}
Simplify $\dfrac{x^2+7x+12}{x+3}$.\workspace{2.4cm}
\end{qbox}

\begin{qbox}{Question 4}
Simplify $\dfrac{x^2-4}{x^2-2x}$.\workspace{2.4cm}
\end{qbox}

\begin{qbox}{Question 5}
Simplify $\dfrac{3x+12}{x^2+4x}$.\workspace{2.4cm}
\end{qbox}

\begin{qbox}{Question 6}
Simplify $\dfrac{x^2-x-6}{x^2-9}$.\workspace{2.4cm}
\end{qbox}

\begin{qbox}{Question 7}
Simplify $\dfrac{x^2-1}{x+2}\cdot\dfrac{x+2}{x-1}$.\workspace{2.4cm}
\end{qbox}

\begin{qbox}{Question 8}
Simplify $\dfrac{x+3}{x}\cdot\dfrac{2x}{x+3}$.\workspace{2.4cm}
\end{qbox}

\begin{qbox}{Question 9}
Simplify $\dfrac{x+2}{x-1}\div\dfrac{x+2}{x}$.\workspace{2.4cm}
\end{qbox}

\begin{qbox}{Question 10}
Simplify $\dfrac{x^2}{x-4}\div\dfrac{x}{x-4}$.\workspace{2.4cm}
\end{qbox}

\begin{qbox}{Question 11}
Simplify $\dfrac{2x^2-8}{x^2-4}$.\workspace{3cm}
\end{qbox}

\begin{qbox}{Question 12}
State the restrictions for $\dfrac{x+1}{x^2-x-6}$.\workspace{2.4cm}
\end{qbox}

\begin{qbox}{Question 13 --- Challenge}
Simplify $\dfrac{x^2-4}{x+3}\cdot\dfrac{x^2+6x+9}{x-2}$.\workspace{3.5cm}
\end{qbox}

\clearpage
\sech{Question Recap \& Answer Key}

\begin{recapbox}
\begin{multicols}{2}
\small
\begin{enumerate}[leftmargin=*,itemsep=3pt]
  \item Simplify $\dfrac{4x^3}{8x}$ and state restrictions.
  \item Simplify $\dfrac{x^2-25}{x+5}$.
  \item Simplify $\dfrac{x^2+7x+12}{x+3}$.
  \item Simplify $\dfrac{x^2-4}{x^2-2x}$.
  \item Simplify $\dfrac{3x+12}{x^2+4x}$.
  \item Simplify $\dfrac{x^2-x-6}{x^2-9}$.
  \item Simplify $\dfrac{x^2-1}{x+2}\cdot\dfrac{x+2}{x-1}$.
  \item Simplify $\dfrac{x+3}{x}\cdot\dfrac{2x}{x+3}$.
  \item Simplify $\dfrac{x+2}{x-1}\div\dfrac{x+2}{x}$.
  \item Simplify $\dfrac{x^2}{x-4}\div\dfrac{x}{x-4}$.
  \item Simplify $\dfrac{2x^2-8}{x^2-4}$.
  \item State the restrictions for $\dfrac{x+1}{x^2-x-6}$.
  \item Simplify $\dfrac{x^2-4}{x+3}\cdot\dfrac{x^2+6x+9}{x-2}$.
\end{enumerate}
\end{multicols}
\end{recapbox}

\begin{keybox}
\begin{multicols}{2}
\small
\begin{enumerate}[leftmargin=*,itemsep=3pt]
  \item $\dfrac{x^2}{2},\ x\ne0$
  \item $x-5,\ x\ne-5$
  \item $x+4,\ x\ne-3$
  \item $\dfrac{x+2}{x},\ x\ne0,2$
  \item $\dfrac{3}{x},\ x\ne0,-4$
  \item $\dfrac{x+2}{x+3},\ x\ne\pm3$
  \item $x+1,\ x\ne-2,1$
  \item $2,\ x\ne0,-3$
  \item $\dfrac{x}{x-1},\ x\ne0,1,-2$
  \item $x,\ x\ne0,4$
  \item $2,\ x\ne\pm2$
  \item $x\ne3,-2$
  \item $(x+2)(x+3),\ x\ne-3,2$
\end{enumerate}
\end{multicols}
\end{keybox}

\end{document}
